\documentclass[12pt]{letter}
\usepackage[utf8]{inputenc}
\usepackage[T1]{fontenc}
\usepackage{geometry}
\geometry{margin=1in}
\usepackage{hyperref}

\signature{Ian Todd\\Sydney Medical School\\University of Sydney}
\address{Ian Todd\\Sydney Medical School\\University of Sydney\\Sydney, NSW, Australia\\itod2305@uni.sydney.edu.au}

\begin{document}

\begin{letter}{Professor A.U.\ Igamberdiev\\Editor, BioSystems}

\opening{Dear Professor Igamberdiev,}

Please find enclosed the revised manuscript BIOSYS-D-25-00981, ``Coherence time in biological oscillator assemblies bounds the rate of state registration,'' for consideration as a revised submission to \textit{BioSystems}.

I thank both reviewers for their constructive feedback. All points raised have been addressed, and the revision includes substantial additional developments.

\medskip
\textbf{Response to Reviewer \#1}

\medskip
\textit{Reviewer:} ``The speed of thinking is restricted by the Landauer and Margolus-Levitin bounds\ldots\ This should be discussed in the revised version.''

\textit{Response:} Section 2.4 (``Quantum and thermodynamic limits of computation'') now explicitly addresses the Landauer bound, the Margolus-Levitin quantum speed limit, and the energy-time uncertainty relation. For neural-scale energies, QSL and Landauer bounds are far below observed cognitive timescales (by approximately eleven orders of magnitude) and are therefore not rate-limiting in this regime. The power limit ($\tau_{\mathrm{power}}$) is now quantified with a detailed back-of-envelope calculation (\S2.1) showing it is non-binding even at biophysical (not merely Landauer-floor) energy scales.

\medskip
\textbf{Response to Reviewer \#2}

\medskip
\textit{Reviewer:} ``Speed of thought should be changed to speed of neural processes associated with thinking.''

\textit{Response:} ``Speed of thought'' has been replaced with ``neural processes associated with cognition'' throughout, including the abstract.

\medskip
\textit{Reviewer:} ``Quantum limits of information transmission based on energy-time uncertainty relation should be discussed.''

\textit{Response:} Energy-time uncertainty is discussed in \S2.4, following the framework developed by Igamberdiev (2021, 2025, 2026).

\medskip
\textit{Reviewer:} Cite four specific BioSystems papers (DOIs provided).

\textit{Response:} All four are now cited: Igamberdiev (2021, 104471), Igamberdiev (2025, 105451), Igamberdiev (2026, 105707), and Khrennikov et al.\ (2025, 105573).

\medskip
\textit{Reviewer:} ``Abstract should not contain references. Also, if formulas cannot be avoided in abstract, it would be better that you simplify them wherever possible.''

\textit{Response:} The abstract contains no citation numbers. Mathematical notation has been simplified; only simple inline symbols ($M$, $r$, $\varepsilon$) are retained for precision.

\medskip
\textbf{Additional improvements}

\medskip
Beyond the reviewer requests, the revision includes several substantial developments:

\begin{itemize}
\item \textbf{Phase delta regime} (\S2.3): New section identifying the pre-commit interval---in which structured inter-module phase relationships bias downstream outcomes without discrete state registration---as a computationally productive mode operating below the Landauer threshold. Includes formal Definition~2 and the speed-flexibility trade-off (Proposition~2).

\item \textbf{Substrate identification} (\S2.1): Argument from the framework's own definitions that the extracellular electric field is the physical carrier of the phase delta regime: it is the only substrate in neural tissue that is simultaneously continuous, spatially extended, and high-dimensional. The causal significance of ephaptic coupling follows from two premises: fields modulate spike timing (Anastassiou et al.\ 2011), and timing is causally significant (the paper's central result). The magnitude of the effect scales with effective dimensionality $D_{\mathrm{eff}}$.

\item \textbf{Power limit quantified} (\S2.1): Back-of-envelope calculation showing $\tau_{\mathrm{power}} \sim 10^{-15}$--$10^{-8}$~s, non-binding even at biophysical energy scales. Added to the hierarchy listing in \S2.4.

\item \textbf{Improved validation methodology}: Kaplan-Meier survival analysis for censored first-passage times, theory-consistent fitting ($\hat{\alpha} = 0.99$, $R^2 = 0.97$), and regime-boundary diagnostics across three network topologies.

\item \textbf{New sections}: Frequency-dimensionality relationship (\S4.5), subjective temporal scaling (\S4.6) with pharmacological predictions across four drug classes from a two-parameter structure.

\item \textbf{Alpha oscillations and perception sets} (\S4.4): Integrated recent work by Wutz (2024) showing alpha reflects internally generated ``perception sets.'' We show this is complementary: a perception set is a high-dimensional internal template, and the temporal effects documented---dilation under richer processing, compression under impoverished conditions---correspond to the exponential dependence of $\tau_{\mathrm{coh}}$ on coordination depth.

\item \textbf{Synchronisation theory connections} (\S5): Added connections to the EI-Kuramoto model (Kuroki \& Mizuseki, 2025) and criticality as a computational operating point (Hengen \& Shew, 2025). Added Buzs\'{a}ki \& Draguhn (2004) for the frequency-hierarchy view.

\item \textbf{Expanded limitations}: Parameter identifiability, regime boundaries ($R^2 = 0.28$ for non-modular networks), module-size confound analysis, dwell-time sensitivity.
\end{itemize}

The manuscript has grown from 41 to 55 pages to accommodate these developments.

\medskip
\textbf{Data and code availability.} All simulation code and parameter estimation workflows are available at \url{https://github.com/todd866/coherence-time}.

\closing{Sincerely,}

\end{letter}
\end{document}
